\chapter{Client Pain Points}

Rather than repeat the sector-by-sector breakdown given in the sector analysis,
this chapter synthesises the pain points that recur across sectors and isolates
the single constraint that determines whether any of them justifies a robot in
the Tunisian context.

\section{Recurring Operational Pain Points}

Across factories, warehouses, logistics hubs, hospitals, hotels, universities,
and security operations, the operational complaints cluster into a small number
of recurring themes:

\begin{itemize}
    \item \textbf{Manual internal transport.} Moving parts between workstations
          (factories), pallets and orders (warehouses, logistics), medication and
          samples (hospitals), room deliveries (hotels), and mail between
          buildings (universities) is slow, repetitive, and fatigue-prone. In
          hospitals specifically, clinical staff can spend a large share of their
          time on transport tasks rather than patient care
          \cite{researchgate_hospital,mir_healthcare}.
    \item \textbf{Repetitive and hazardous inspection rounds.} Manual gauge
          reading, thermal and leak detection, and equipment monitoring are slow,
          infrequent, and dangerous in some zones, with faults often detected
          late.
    \item \textbf{Labor scarcity and demand peaks.} Seasonal volume swings
          (logistics, tourism) and turnover in repetitive roles are hard to
          absorb with temporary staff.
    \item \textbf{Data and coverage gaps.} Inconsistent inspection schedules and
          fixed-camera blind spots leave facilities without a continuous record
          of their own condition.
\end{itemize}

\section{The Overarching Constraint: Cost-to-Labor Ratio}

The pain points above are largely shared with the high-wage markets that
existing AMRs were designed for. What changes locally is not the pain but the
economics of solving it. Because Tunisian wages are substantially lower than in
Europe or the GCC, a robot bought to replace a transport worker has a much
longer payback than the same robot in a German factory --- a point international
product sheets never make, because they are written for high-labor-cost markets.

The practical consequence is that pure labor substitution is rarely the pain
worth solving here. The cases where the economics work are those where the robot
is compared against something other than a wage:

\begin{itemize}
    \item \textbf{Inspection}, where the robot is measured against the cost of an
          avoided incident, downtime, or undetected failure --- which can exceed
          the cost of a simple inspection robot even for a Tunisian factory.
    \item \textbf{Universities}, where the robot is measured against the cost of
          an imported research platform that institutions already buy, not
          against a salary.
\end{itemize}

This distinction is the connective thread running through every sector
assessment in this report.

\section{Financial and Procurement Pain}

A second, cross-cutting layer of pain is financial rather than operational.
Local buyers are capex-averse and procurement is slow, particularly in the
public sector. Imported hardware adds foreign-currency exposure and customs and
lead-time friction. Thin-margin operators (3PL, retail) cannot absorb a large
capital outlay at all. These constraints make an outright-purchase model a poor
fit and push strongly toward a service-based commercial model
\cite{builtin_raas}, which is addressed in the pricing chapter.
